\subsection{Relations between the properties of a Noetherian local
ring \texorpdfstring{$A$}{A} and of a quotient ring
\texorpdfstring{$A/tA$}{A/tA}}
\label{subsection:IV.5.12}
\label{IV.5.12}

\oldpage[IV-2]{126}

We have already seen in \sref{IV.3.4} relations between the
properties of $A$ and those of $A/tA$ concerning associated prime
ideals, as well as the properties of being a domain or being
reduced.  In this subsection we give further relations between the
properties of these rings, involving the notions of dimension and
depth.

\begin{proposition}[5.12.1]
\label{IV.5.12.1}
Let $A$ be a Noetherian local ring, put $X=\Spec(A)$, and let $t$
be an element of $A$ belonging to a system of parameters
\sref[0]{0.16.3.6}.  Let $X_0$ be the closed subspace $V(t)$ of
$X$, and let $X_1=X-X_0$.  Let $Z$ be a subset of $X$ stable
under specialization.  Suppose that $X$ is biequidimensional
(as is the case when $A$ is equidimensional and is a quotient of
a regular local ring \sref[0]{0.16.5.12}).  If
$Z_0=Z\cap X_0$ and $Z_1=Z\cap X_1$, then
\begin{equation}
\tag{5.12.1.1}
\label{IV.5.12.1.1}
  \codim(Z_0,X_0)\leq\codim(Z,X)\leq\codim(Z_1,X_1).
\end{equation}
\end{proposition}

\begin{proof}
The second inequality follows from the definition
\sref{IV.5.1.3}.  For the first, it is enough to treat the case
in which $Z$ is closed.  Indeed, by hypothesis $Z$ is a union of
closed subsets $Z_\alpha$, and inequalities
\[
  \codim(X_0\cap Z_\alpha,X_0)\leq\codim(Z_\alpha,X)
\]
for all $\alpha$ give
\[
\begin{split}
  \codim(Z_0,X_0)
  &=\inf_\alpha\codim(X_0\cap Z_\alpha,X_0)\\
  &\leq\inf_\alpha\codim(Z_\alpha,X)
   =\codim(Z,X).
\end{split}
\]
Suppose therefore that $Z$ is closed.  Then
\[
  \codim(Z,X)=\dim(X)-\dim(Z)
\]
by \sref[0]{0.14.3.5.1}.

\oldpage[IV-2]{127}

Moreover, $X_0$ is catenary and equidimensional, and
$\dim(X_0)=\dim(X)-1$ by \sref[0]{0.16.3.4}; hence
\[
  \codim(Z_0,X_0)
  =\dim(X)-1-\dim(Z_0).
\]
Since $\dim(Z_0)\geq\dim(Z)-1$ \sref[0]{0.16.3.4}, the first
inequality in \eqref{IV.5.12.1.1} follows.
\end{proof}

\begin{proposition}[5.12.2]
\label{IV.5.12.2}
Let $A$ be a catenary Noetherian local ring, let $M$ be a finite
$A$-module, let $t$ be an $M$-regular element belonging to the
maximal ideal $\mathfrak m$ of $A$, and let $k\geq1$.  If $M/tM$
is equidimensional and satisfies $(\mathrm S_k)$, then the same is
true of $M$.
\end{proposition}

\begin{proof}
In view of $(\mathbf{Err}_{\mathrm{III}},30)$ and
\sref{IV.5.7.3}[vi], we may reduce to the case
\[
  \operatorname{Supp}(M)=\Spec(A)=X.
\]
Put $X_0=V(t)=\Spec(A/tA)$.  Since $M/tM$ satisfies
$(\mathrm S_1)$, it has no embedded associated prime ideals
\sref{IV.5.7.5}.  Applying \sref{IV.3.4.4} at the closed point of
$X$ (and of every closed subprescheme of $X$), we see that the
irreducible components of $\operatorname{Supp}(M/tM)$ are exactly
the irreducible components of $Y_i\cap X_0$, where
$Y_i$ $(1\leq i\leq r)$ are the irreducible components of $X$.

Because $t$ is $M$-regular, $V(t)$ contains no maximal point of
$X$ \sref{IV.3.1.8}.  Every irreducible component of
$Y_i\cap X_0$ therefore has codimension $1$ in $Y_i$
\sref{IV.5.1.8}.  Since $X$ is catenary and all irreducible
components of $\operatorname{Supp}(M/tM)$ have the same
dimension, the $Y_i$ all have the same dimension.  Thus $X$ is
equidimensional, and hence biequidimensional because $A$ is local
and catenary.

It remains to prove $(\mathrm S_k)$.  By the criterion
\sref{IV.5.7.4}, with its notation, we must show that
\[
  \codim(Z_n,X)>n+k
\]
for every integer $n\geq0$.  The hypothesis on $M/tM$ and the same
criterion give
\[
  \codim(Z_n\cap X_0,X_0)>n+k.
\]
Every specialization of a point of $Z_n$ again belongs to $Z_n$,
as will be proved in \sref{IV.6.11.5}.\footnote{The reader may
verify that {\hyperlink{ega4.c04.prop.5-12-2}{Proposition~5.12.2}} is not used in the proof of
\sref{IV.6.11.5}.}  Finally, $t$ belongs to a system of
parameters for $M$ \sref[0]{0.16.4.1}, and hence also to one for
$A$, because the annihilator of $M$ is nilpotent under the
hypothesis $\operatorname{Supp}(M)=\Spec(A)$.  The required
inequality now follows from \sref{IV.5.12.1}.
\end{proof}

\begin{remark}[5.12.3]
\label{IV.5.12.3}
If in {\hyperlink{ega4.c04.prop.5-12-2}{Proposition~5.12.2}} one assumes only that $M/tM$ is
equidimensional, it need not follow that $M$ is
equidimensional.  For example, let $k$ be a field, let
$B=k[T,U,V]$, let $C$ be the local ring of $B$ at the maximal
ideal $BT+BU+BV$, and put
\[
  \mathfrak p=CU,\qquad
  \mathfrak q=CV+C(T+U),\qquad
  A=C/\mathfrak p\mathfrak q.
\]
Geometrically, $A$ is the local ring, at their intersection
point $x$, of the closed subscheme of affine $3$-space formed by
a plane and a line meeting that plane at $x$.  Let $t,u,v$ be the
images of $T,U,V$ in $A$.  The element $t$ is not a zero divisor
in $A$, and
\[
  A/tA\simeq C_0/\mathfrak p_0\mathfrak q_0,
\]
where $C_0$ is the local ring of $B_0=k[U,V]$ at the maximal
ideal $B_0U+B_0V$, while
\[
  \mathfrak p_0=C_0U,\qquad
  \mathfrak q_0=C_0U+C_0V.
\]
Thus $\mathfrak q_0$ is the maximal ideal of $C_0$.  It is
immediate that $\Spec(A/tA)$ is irreducible of dimension $1$,
that $A/tA$ does not satisfy $(\mathrm S_1)$, and that
$\Spec(A)$ is not equidimensional.
\end{remark}

\begin{corollary}[5.12.4]
\label{IV.5.12.4}
Let $A$ be a catenary Noetherian local ring, let $t$ be a regular
element belonging to the maximal ideal of $A$, and let $k\geq1$.
If $A/tA$ satisfies $(\mathrm S_k)$, then so does $A$.
\end{corollary}

\begin{proof}
For $k=1$ this follows from \sref{IV.3.4.4}, in view of the
interpretation of $(\mathrm S_1)$ in \sref{IV.5.7.5}.  For
$k\geq2$, Hartshorne's criterion \sref{IV.5.10.9} shows that
$A/tA$ is equidimensional, and {\hyperlink{ega4.c04.prop.5-12-2}{Proposition~5.12.2}} applies.
\end{proof}

\begin{proposition}[5.12.5]
\label{IV.5.12.5}
Let $A$ be a catenary Noetherian local ring, let $t$ be a regular
element of $A$ belonging to its maximal ideal, and let $k\geq0$.
If $A/tA$ is reduced and equidimensional and satisfies
$(\mathrm R_k)$, then $A$ has these same three properties.
\end{proposition}

\begin{proof}
We already know from \sref{IV.3.4.6} that $A$ is reduced, and
{\hyperlink{ega4.c04.prop.5-12-2}{Proposition~5.12.2}}, applied with $k=1$, shows that $A$ is
equidimensional.  Put
\[
  X=\Spec(A),\qquad X_0=V(t)=\Spec(A/tA),
\]
and let $Z$ be the set of points at which $X$ is not regular.
Every specialization of a point of $Z$ belongs to $Z$
\sref[0]{0.17.3.2}.  On the other hand, if $x\in X_0$ and $X_0$
is regular at $x$, then $X$ is regular at $x$
\sref[0]{0.17.1.8}.  Consequently the nonregular locus $Z'_0$
of $X_0$ contains $Z_0=Z\cap X_0$.  The hypothesis gives
\[
  k<\codim(Z'_0,X_0)\leq\codim(Z_0,X_0).
\]
Since $A$ is equidimensional and catenary, {\hyperlink{ega4.c04.prop.5-12-1}{Proposition~5.12.1}}
gives
\[
  \codim(Z_0,X_0)\leq\codim(Z,X).
\]
Thus $X$ satisfies $(\mathrm R_k)$.
\end{proof}

\begin{remark}[5.12.6]
\label{IV.5.12.6}
If in {\hyperlink{ega4.c04.prop.5-12-5}{Proposition~5.12.5}} one assumes only that $A/tA$ is reduced
and satisfies $(\mathrm R_k)$, it need not follow that $A$
satisfies $(\mathrm R_k)$.  For example, let $k$ be a field and
let $P(U,V)\in k[U,V]$ be irreducible, with the curve $\Gamma$
defined by $(P)$ in $\Spec(k[U,V])$ singular at the maximal ideal
$(U)+(V)$.  One may take, for example,
\[
  P(U,V)=U(U^2+V^2)+(U^2-V^2),
\]
which defines a cubic with a double point.  Let
$B=k[T,U,V,W]$, let $C$ be the local ring of $B$ at the maximal
ideal $BT+BU+BV+BW$, and put
\[
  \mathfrak p=CW+CP(U-T,V),\qquad \mathfrak q=CU.
\]
Consider the local ring $A=C/\mathfrak p\mathfrak q$.
Geometrically, the corresponding closed subscheme $X$ of affine
$4$-space is the union of a hyperplane $H$ and a
$2$-dimensional ``cylinder'' $L$ with base $\Gamma$; the
``singular line'' $Y$ of $L$ is not contained in $H$, and $A$ is
the local ring of $X$ at the point $x=Y\cap H$.

Let $t$ be the image of $T$ in $A$.  The scheme
$\Spec(A/tA)$ has $x$ as its only singular point; its local ring
there is $A/tA$, which is reduced and has dimension $2$.  By
contrast, the generic point $y$ of $Y$, defined by the image in
$A$ of
\[
  C(U-T)+CV+CW,
\]
is a singular point of $X$, and $\sh O_{X,y}$ has dimension $1$.
Thus $A/tA$ is reduced and satisfies $(\mathrm R_1)$, whereas
$A$ does not satisfy $(\mathrm R_1)$.
\end{remark}

\oldpage[IV-2]{129}

\begin{corollary}[5.12.7]
\label{IV.5.12.7}
Let $A$ be a catenary Noetherian local ring and let $t$ be a
regular element of $A$ belonging to its maximal ideal.  If
$A/tA$ is a domain and is integrally closed, then so is $A$.
\end{corollary}

\begin{proof}
By Serre's criterion \sref{IV.5.8.6} and the fact that $A$ is
local, it is enough to prove that $\Spec(A)$ satisfies
$(\mathrm S_2)$ and $(\mathrm R_1)$.  The hypothesis on $A/tA$
implies $(\mathrm R_1)$ for $A$ by {\hyperlink{ega4.c04.prop.5-12-5}{Proposition~5.12.5}} and
$(\mathrm S_2)$ for $A$ by {\hyperlink{ega4.c04.cor.5-12-4}{Corollary~5.12.4}}.
\end{proof}
